\chapter{How to write a proof}\label{ch:proof}
\index{proof}
% chapter-local: amber defect marker used in Figure~\ref{fig:markup} and text
\newcommand{\defect}[1]{\textcolor{amber}{\sffamily\bfseries\textsuperscript{[#1]}}}

Chapter~\ref{ch:logic} gave you the grammar claims are written in. This chapter
gives you the machinery that establishes them: a small inventory of proof
techniques --- direct, cases, contrapositive, contradiction, two directions for
an iff, witnesses for existence, counterexamples for disproof --- and, just as
important, the writing habits that make a proof readable and checkable. After
this chapter you will be able to look at a claim, name the technique its shape
asks for, and produce a proof a grader can verify line by line. Every later
chapter is practice on this one; the techniques here never change, only the
raw material does.

\section{What a proof is}\label{sec:whatproof}
\index{axiom}\index{inference}

\begin{definition}[Proof]\index{proof!formal definition}
A \emph{proof} of a statement is a finite chain of assertions ending at that
statement, where each assertion is either an axiom, a definition, a hypothesis,
or a consequence of earlier lines by a rule of inference.
\end{definition}

Three kinds of raw material can appear in the chain.
\emph{Axioms}\index{axiom} are the ground rules you do not prove: for this book,
the arithmetic laws of $\Z$, $\Q$ and $\R$ (addition, multiplication, order),
plus the definitions Chapter~\ref{ch:logic} fixed. \emph{Definitions} unfold on
demand: ``even'' always expands to ``$=2k$ for some integer $k$''.
\emph{Theorems}\index{theorem} are claims already proved; once proved, a theorem
becomes a tool you may cite by name instead of re-proving it.

The legal moves between lines are the rules of \emph{inference}, which are just
Chapter~\ref{ch:logic} applied: from $P$ and $P \Rightarrow Q$ you may write $Q$
(\emph{modus ponens}\index{modus ponens}); from $\forall x\,P(x)$ you may write
$P(a)$ for a specific $a$; from $P \land Q$ you may write $P$. If a line of your
proof cannot name its licence --- ``by definition of even'', ``by modus ponens on
lines 2 and 4'' --- the line does not belong yet.

\begin{remark}
``Axioms vs.\ proven tools'' is worth keeping straight for one practical reason:
in an exam you may cite theorems you have proved or been given, but not the
claim you are currently proving. Citing your own goal is the circularity trap
of Section~\ref{sec:mistakes2}.
\end{remark}

\section{Direct proof and the template}\label{sec:direct}
\index{direct proof}

The workhorse. To prove $P \Rightarrow Q$: assume $P$, derive $Q$. Almost every
proof you write this year will carry this skeleton, often with other techniques
nested inside.

\begin{principle}[The assume--derive--conclude template]
\index{proof!template}
Write the first and last lines before the middle:
\emph{Assume} the hypothesis (unfolded to definitions), \emph{conclude} the goal
(restated), then \emph{derive} the gap. Figure~\ref{fig:prooftemplate} shows the
shape; the middle is where the problem lives, but the two ends are free.
\end{principle}

\begin{figure}[ht]\centering
% Figure: the direct-proof template --- assume, derive, conclude.
% Source for Figure~\ref{fig:prooftemplate}. TikZ only; regenerable by compiling main.tex.
\begin{tikzpicture}[
  stage/.style={draw,rounded corners=2pt,inner sep=6pt,font=\small,text width=64mm,align=flush left},
  arr/.style={->,accent,thick},
  lab/.style={font=\footnotesize\sffamily,text=slate,text width=36mm,align=flush left},
]
\node[stage] (a) at (0,0) {\textbf{Assume.} Unfold the hypotheses into their
definitions --- ``$n$ is even'' becomes ``$n = 2k$ for some integer $k$''.};
\node[stage] (d) at (0,-2.3) {\textbf{Derive.} Write a chain of steps; each line
is licensed by a definition, an axiom, or a theorem already proved.};
\node[stage] (c) at (0,-4.5) {\textbf{Conclude.} Stop the instant the goal lands
and restate it --- ``therefore $n+m$ is even''.};
\draw[arr] (a.south) -- (d.north);
\draw[arr] (d.south) -- (c.north);
\node[lab,anchor=north west] at (3.7,0.1) {the hypotheses are the only raw
material you start with};
\node[lab,anchor=north west] at (3.7,-2.25) {every step cites its licence ---
nothing enters ``obviously''};
\node[lab,anchor=north west] at (3.7,-4.4) {one clean landing; do not keep
writing past it};
\end{tikzpicture}

\caption{The direct-proof template. Notice that the creative part is confined to
the middle stage --- the assume line and the conclude line are determined by the
claim itself, so writing them down first costs nothing and fixes the target.}
\label{fig:prooftemplate}
\end{figure}

\begin{theorem}\label{thm:evensum}
If $m$ and $n$ are even integers, then $m + n$ is even.
\end{theorem}

\begin{proof}
Let $m$ and $n$ be even integers. By definition of even, $m = 2j$ and $n = 2k$
for some integers $j$ and $k$. Then $m + n = 2j + 2k = 2(j + k)$, and $j + k$ is
an integer, so $m + n$ is even.
\end{proof}

Read the proof again against the template: the assume line unfolds ``even'' on
both inputs at once; the derivation is a single algebraic rewrite; the conclude
line checks the definition's existential (``for some integer'') by exhibiting
$j+k$. That last detail --- the new object $j+k$ must be shown to be an integer
--- is what separates a proof from a gesture.

\begin{remark}
\index{fresh names}
Two even numbers need two names: $m = 2j$, $n = 2k$. Writing $m = 2k$ and
$n = 2k$ forces $m = n$, a silent extra assumption. Fresh objects get fresh
names --- one of the cheapest correctness rules in this book.
\end{remark}

\begin{exercise}[easy]\label{exr:oddsum}
Prove that the sum of two odd integers is even, using the template. (Odd means
$2k+1$ for some integer $k$.)
\end{exercise}

\section{Proof by cases}\label{sec:cases}
\index{proof!by cases}

Sometimes the hypothesis does not unfold into one usable form, but the universe
splits into a few shapes that each do. \emph{Proof by cases} handles every shape
separately; it is valid exactly when the cases are \emph{exhaustive} --- every
object falls into at least one case.

\begin{theorem}\label{thm:n2n}
For every integer $n$, the number $n^2 + n$ is even.
\end{theorem}

\begin{proof}
Every integer is even or odd (Section~\ref{sec:parity}), so the two cases cover
all $n$.
\emph{Case 1: $n$ even.} Write $n = 2k$. Then $n^2 + n = 4k^2 + 2k = 2(2k^2 + k)$,
and $2k^2 + k$ is an integer, so $n^2 + n$ is even.
\emph{Case 2: $n$ odd.} Write $n = 2k + 1$. Then
$n^2 + n = 4k^2 + 6k + 2 = 2(2k^2 + 3k + 1)$, again even.
\end{proof}

\begin{wins}[Cases done right]
\index{cases!exhaustiveness}
Before the cases, \textbf{say why they are exhaustive} (``every integer is even
or odd''). A third case you forgot is not a gap in style --- it is a hole in the
proof. When the split is by residues ($n \bmod 4 \in \{0,1,2,3\}$), exhaustiveness
is automatic; when it is by geometry or sign, state it explicitly.
\end{wins}

\begin{exercise}[medium]\label{exr:three}
Prove that for every integer $n$, the product $n(n+1)(n+2)$ is divisible by $6$.
(Cases on $n \bmod 6$? Or: is one factor always divisible by $2$, and one by $3$?)
For the contest version of divisibility-by-casework, see
AMC~12~2023~A~P3~\cite{aops-amc}.
\end{exercise}

\section{Contrapositive}\label{sec:contrapositive}
\index{contrapositive!proof by}

Chapter~\ref{ch:logic} showed $P \Rightarrow Q$ and $\lnot Q \Rightarrow \lnot P$
carry the same truth value (Figure~\ref{fig:family}). So they are the same job:
prove whichever unfolds more usefully. The contrapositive earns its keep exactly
when the \emph{hypothesis} has no handle but the \emph{negated conclusion} does.

\begin{example}\label{ex:n2even}
Prove: for every integer $n$, if $n^2$ is even then $n$ is even.
\par\smallskip
\textbf{Strategy.} Try it directly first and watch where it stalls: $n^2$ even
gives $n^2 = 2k$, and now you must say something about $n$ --- but $n$ is
$\sqrt{2k}$, and square roots of integers need not be integers. The hypothesis
talks about $n^2$; nothing in it unfolds $n$. Flip to the contrapositive, where
the hypothesis is ``$n$ is odd'' --- a form of $n$ itself, which is exactly the
handle we lacked.
\par\smallskip
\begin{solution}
We prove the contrapositive: if $n$ is odd, then $n^2$ is odd. Let $n$ be odd,
so $n = 2k + 1$ for some integer $k$. Then
\[
n^2 = (2k+1)^2 = 4k^2 + 4k + 1 = 2(2k^2 + 2k) + 1,
\]
and $2k^2 + 2k$ is an integer, so $n^2$ is odd. By the contrapositive, if $n^2$
is even then $n$ is even.
\end{solution}
\end{example}

Two habits to copy from Example~\ref{ex:n2even}: the proof \emph{says out loud}
that it is proving the contrapositive (a reader cannot guess your route), and
the strategy block shows the stalled direct attempt --- knowing why the obvious
route fails is how you learn to pick the right one.

\begin{exercise}[easy]\label{exr:sumge2}
Prove by contrapositive: for real numbers $x, y$, if $x + y \ge 2$ then
$x \ge 1$ or $y \ge 1$.
\end{exercise}

\section{Contradiction}\label{sec:contradiction}
\index{proof!by contradiction}

To prove $P$ by contradiction: assume $\lnot P$ and derive an impossibility ---
something colliding with a definition, an axiom, or arithmetic itself. It is the
technique of last resort and of necessity: for claims of the form ``no object
satisfies\dots'' it is often the only door, since a universal negative gives you
nothing to construct.

\begin{principle}[Contradiction template]
\index{contradiction}
``Assume, for contradiction, that \emph{the claim fails}.'' Then treat the
negated claim --- usually $\exists$ of something forbidden --- as a new object
in hand and squeeze it until it breaks.
\end{principle}

\begin{example}[$\sqrt{2}$ is irrational]\label{ex:sqrt2}
Prove: there are no integers $p, q$ with $q \ne 0$ and $(p/q)^2 = 2$.
\par\smallskip
\textbf{Strategy.} The claim is a non-existence claim --- the natural habitat of
contradiction. Suppose a fraction $p/q$ in lowest terms squares to $2$ (``lowest
terms'' is the extra lever: it makes the witness \emph{minimal}, and minimal
objects can be forced below themselves). Unfolding $p^2 = 2q^2$ shows $p$ is
even; that forces $q$ even too; and two even numbers were never in lowest terms.
\par\smallskip
\begin{solution}
Assume, for contradiction, that $\sqrt{2} = p/q$ for integers $p, q$ with
$q \ne 0$, written in lowest terms, so $p$ and $q$ share no common factor.
Then $p^2 = 2q^2$, so $p^2$ is even, so by Example~\ref{ex:n2even} $p$ is even:
write $p = 2a$. Substituting, $4a^2 = 2q^2$, hence $q^2 = 2a^2$ is even, so $q$
is even too. Then $p$ and $q$ share the factor $2$, contradicting lowest terms.
Therefore no such $p, q$ exist: $\sqrt{2}$ is irrational.
\end{solution}
\end{example}

\begin{example}[Working modulo $4$]\label{ex:mod4}
Prove: there are no integers $a, b$ with $a^2 - 4b = 2$.
\par\smallskip
\textbf{Strategy.} Again a non-existence claim, so contradiction. But squeeze
\emph{where}? The term $4b$ is invisible modulo $4$: reducing the equation mod
$4$ erases $b$ entirely and leaves the finite question ``can $a^2$ be $2$ mod
$4$?'' --- checkable by cases on $a$. Modular arithmetic is a lens: it discards
the parts of a statement that cannot obstruct you and keeps the part that can.
\par\smallskip
\begin{solution}
Assume integers $a,b$ satisfy $a^2 - 4b = 2$. Reducing mod $4$ gives
$a^2 \equiv 2 \pmod{4}$. But $a$ is even or odd: if $a = 2k$ then
$a^2 = 4k^2 \equiv 0$; if $a = 2k+1$ then $a^2 = 4k^2 + 4k + 1 \equiv 1$. So
$a^2 \bmod 4 \in \{0,1\}$, never $2$ --- contradiction. No such $a,b$ exist.
\end{solution}
\end{example}

Notice the proof inside Example~\ref{ex:mod4} is a two-case argument nested in
a contradiction: techniques compose. The mod-$4$ lens will return throughout the
number theory chapters; remainders problems in this style appear constantly on
contests (see AMC~12~2024~B~P14 on the AoPS archive~\cite{aops-amc}).

\begin{exercise}[medium]\label{exr:sqrt3}
Adapt Example~\ref{ex:sqrt2} to prove $\sqrt{3}$ is irrational. (Where does the
argument use that $2$ is prime? What goes wrong for $\sqrt{4}$?)
\end{exercise}

\section{If and only if}\label{sec:iff}
\index{iff!proof of}

A biconditional $P \iff Q$ is two jobs stapled together, and you must do both:
the forward direction $P \Rightarrow Q$ and the reverse $Q \Rightarrow P$ ---
each by whichever technique suits it. Label them ($\Rightarrow$) and
($\Leftarrow$) so the reader can audit both.

\begin{theorem}\label{thm:neveniff}
For every integer $n$: $n$ is even if and only if $n^2$ is even.
\end{theorem}

\begin{proof}
($\Rightarrow$) If $n = 2k$, then $n^2 = 4k^2 = 2(2k^2)$, which is even.
($\Leftarrow$) This is Example~\ref{ex:n2even}, proved there by contrapositive.
\end{proof}

\begin{tripwire}
\index{iff!proving one direction}
Proving one direction and stopping is the single most common broken
biconditional. ``$n$ even iff $n^2$ even'' needs \emph{both} halves; a proof of
($\Rightarrow$) alone leaves the claim unproved --- and in contests, an unproved
direction usually hides the interesting half. See Common mistakes.
\end{tripwire}

\begin{figure}[ht]\centering
% Figure: a biconditional is two separately proved one-way streets.
% Source for Figure~\ref{fig:iffproof}. TikZ only; regenerable by compiling main.tex.
\begin{tikzpicture}[
  claim/.style={draw,inner sep=5pt,font=\small,rounded corners=1pt},
  road/.style={->,very thick,accent},
  lab/.style={font=\footnotesize\sffamily,text=slate,align=center}
]
\node[claim] (P) at (0,0) {$P$};
\node[claim] (Q) at (6,0) {$Q$};
\draw[road] (P.30) .. controls (2,1.3) and (4,1.3) .. (Q.150);
\draw[road] (Q.210) .. controls (4,-1.3) and (2,-1.3) .. (P.330);
\node[lab] at (3,1.5) {($\Rightarrow$) prove $P \Rightarrow Q$};
\node[lab] at (3,-1.5) {($\Leftarrow$) prove $Q \Rightarrow P$};
\node[lab] at (3,-2.6) {Each arrow is its own proof --- often by a
\emph{different} technique.};
\end{tikzpicture}

\caption{A biconditional is two one-way streets. Notice: the two arrows are
\emph{separate} proofs --- each may need a different technique, and finishing
one leaves the claim half-proved.}
\label{fig:iffproof}
\end{figure}

\section{Existence and uniqueness}\label{sec:existence}
\index{existence}\index{uniqueness}

An existential claim $\exists x\,P(x)$ has the cheapest proof of all: exhibit
one $x$ and check it. This is a \emph{constructive}\index{constructive proof}
proof --- you hand over the object. To prove \emph{unique} existence
$\exists!x\,P(x)$, add a second act: any two candidates $x_1, x_2$ with
$P(x_1), P(x_2)$ must satisfy $x_1 = x_2$ (existence $+$ uniqueness, the two
halves from Chapter~\ref{ch:logic}).

Not every existence proof constructs. The classic counterexample: there exist
irrational numbers $a, b$ such that $a^b$ is rational. Consider
$\sqrt{2}^{\sqrt{2}}$: if it is rational, take $a = b = \sqrt{2}$. If it is
irrational, take $a = \sqrt{2}^{\sqrt{2}}$ and $b = \sqrt{2}$ --- then
$a^b = (\sqrt{2}^{\sqrt{2}})^{\sqrt{2}} = \sqrt{2}^2 = 2$, rational. Either way
some pair works; the proof never says which, and it is still a complete proof.
\emph{Non-constructive}\index{non-constructive proof} existence is legal ---
just know which kind you are writing, since only constructive proofs produce
objects you can compute with.

\begin{example}[A unique solution]\label{ex:unique156}
Prove there is a unique positive integer $n$ with $n^2 + n = 156$.
\par\smallskip
\textbf{Strategy.} Two acts again: find one $n$ (search small values ---
$f(n) = n^2+n$ grows fast), then show no other positive integer works. For
uniqueness, monotonicity is the cheapest tool: if $f$ is strictly increasing
on positive integers, two different inputs can never land on $156$.
\par\smallskip
\begin{solution}
\emph{Existence.} $n = 12$: $12^2 + 12 = 144 + 12 = 156$.
\emph{Uniqueness.} Suppose $m$ is a positive integer with $m^2 + m = 156$.
If $m > 12$ then $m^2 + m \ge 13^2 + 13 = 182 > 156$; if $m < 12$ then
$m^2 + m \le 11^2 + 11 = 132 < 156$. So $m = 12$. $\qed$
\end{solution}
\end{example}

\begin{figure}[ht]\centering
% Figure: existence and uniqueness are two separate acts.
% Source for Figure~\ref{fig:uniq}. TikZ only; regenerable by compiling main.tex.
\begin{tikzpicture}[
  panel/.style={draw,rounded corners=2pt,inner sep=8pt},
  obj/.style={circle,draw,inner sep=2pt,minimum size=4.5mm,font=\footnotesize},
  good/.style={circle,draw,fill=accent!25,inner sep=2pt,minimum size=4.5mm,font=\footnotesize},
  lab/.style={font=\footnotesize\sffamily,text=slate,align=center}
]
% left panel: existence
\begin{scope}
\node[lab] at (1.6,1.0) {Act 1 --- existence};
\node[obj] (a1) at (0,0) {\tiny $x_1$};
\node[good] (a2) at (0.9,0.45) {\tiny $x_2$};
\node[obj] (a3) at (1.8,0.1) {\tiny $x_3$};
\node[obj] (a4) at (2.7,0.5) {\tiny $x_4$};
\node[lab] at (1.6,-0.9) {exhibit \emph{one} that works\\ (the filled node)};
\node[panel] at (1.55,-0.05) {\phantom{XXXXXXX}};
\end{scope}
% right panel: uniqueness
\begin{scope}[xshift=6.4cm]
\node[lab] at (1.6,1.0) {Act 2 --- uniqueness};
\node[obj] (b1) at (0,0.4) {\tiny $x_1$};
\node[obj] (b2) at (1.1,0.4) {\tiny $x_2$};
\draw[->,thick,rulegray] (1.7,0.4) -- (2.4,0.4);
\node[good] (b3) at (3.1,0.4) {\tiny $x$};
\node[lab,align=center]{any two candidates\\ must coincide: $x_1 = x_2$};
\node[panel] at (1.55,-0.05) {\phantom{XXXXXXX}};
\end{scope}
\end{tikzpicture}

\caption{The two acts of $\exists!$. Notice: existence needs only one working
object, but uniqueness needs an argument covering \emph{all} candidates ---
here, squeezing $m$ between $11$ and $13$.}
\label{fig:uniq}
\end{figure}

\begin{exercise}[easy]\label{exr:witness}
Prove $\exists!x \in \Z$ with $3x - 7 = 2$: exhibit the witness, then prove any
two witnesses are equal.
\end{exercise}

\section{Disproof and counterexamples}\label{sec:disproof}
\index{disproof}\index{counterexample}

To refute $\forall x\,P(x)$ you owe one thing only: a single $x$ where $P$
fails. The negation is $\exists x\,\lnot P(x)$, and existentials cost one
witness. Where proofs get long, disproofs get short --- \emph{if} the
counterexample is real. Convention: hunt the \emph{smallest} candidate where the
claim could break, and check it honestly.

\begin{example}[Euler's polynomial]\label{ex:euler}
For which integers $x \ge 0$ is $f(x) = x^2 + x + 41$ prime?
\par\smallskip
\textbf{Strategy.} Evaluate the first values and you get a famous run: $f(0) =
41$, $f(1) = 43$, $f(2) = 47$, all prime --- and the run continues far past
where most people stop checking. The trap is to call the claim proved by
accumulated evidence; the discipline is to ask where it \emph{must} break. Since
$f(x) = x(x+1) + 41$, the value is a multiple of $41$ whenever $41$ divides
$x(x+1)$ --- first at $x = 40$ (and $x = 41$). That is the smallest forced
failure. Check it before believing the streak.
\par\smallskip
\begin{solution}
$f(40) = 40\cdot 41 + 41 = 41^2 = 1681$, composite. So the claim ``$f(x)$ is
prime for all $x \ge 0$'' is false; the counterexample is $x = 40$.
Figure~\ref{fig:euler} and Listing~\ref{lst:euler} show how far the run actually
goes.
\end{solution}
\end{example}

\begin{figure}[ht]\centering
% Figure: Euler's polynomial streak --- f(x)=x^2+x+41 for x=0..50.
% Prime values are lower ticks; composites are marked above with their factors.
% Source for Figure~\ref{fig:euler}. TikZ only; regenerable by compiling main.tex.
\begin{tikzpicture}[
  tick/.style={draw=good,line width=1.1pt},
  bad/.style={draw=amber,line width=1.3pt},
  lab/.style={font=\footnotesize\sffamily,text=slate},
  badlab/.style={font=\footnotesize\sffamily,text=amber,align=center},
  axisline/.style={draw=rulegray},
]
% number line x = 0..50, scaled
\draw[axisline] (0,0) -- (10.6,0);
% decade labels
\foreach \x in {0,10,20,30,40,50}
  \node[lab,anchor=north] at (0.2*\x,-0.32) {\x};
% prime ticks below (all x except composites 40,41,44,49)
\foreach \x in {0,1,2,3,4,5,6,7,8,9,11,12,13,14,15,16,17,18,19,21,22,23,24,25,26,27,28,29,31,32,33,34,35,36,37,38,39,42,43,45,46,47,48,50}
  \draw[tick] (0.2*\x,-0.02) -- (0.2*\x,-0.22);
% composite ticks above the line
\draw[bad] (8.0,0.02) -- (8.0,0.30);
\draw[bad] (8.2,0.02) -- (8.2,0.30);
\draw[bad] (8.8,0.02) -- (8.8,0.30);
\draw[bad] (9.8,0.02) -- (9.8,0.30);
% factorisations as one legend block (ticks carry position, text carries factors)
\node[badlab,anchor=south] at (5.3,0.36)
  {$x{=}40$: $1681 = 41^2$ \qquad $x{=}41$: $1763 = 41\cdot 43$\\
   $x{=}44$: $2021 = 43\cdot 47$ \qquad $x{=}49$: $2491 = 47\cdot 53$};
% streak bracket under the prime run (manual: two drops + a rule)
\draw[rulegray] (0.02,-0.62) -- (0.02,-0.78);
\draw[rulegray] (7.8,-0.62) -- (7.8,-0.78);
\draw[rulegray] (0.02,-0.78) -- (7.8,-0.78)
  node[midway,lab,below=2pt]{40 consecutive primes, $x = 0,\dots,39$};
\end{tikzpicture}

\caption{The streak that is not a theorem. Notice the first composite lands at
the first $x$ with $41 \mid x(x+1)$, exactly where the algebra says it must ---
forty consecutive primes proved nothing on their own.}
\label{fig:euler}
\end{figure}

\begin{checkpoint}[Predict before you run]
Listing~\ref{lst:euler} tests $x = 0$ through $50$. Before running: predict the
first composite $x$, predict how many composites appear in $0..50$, and predict
the smallest factor at $x = 40$ and $x = 41$ without dividing anything ---
use $f(x) = x(x+1) + 41$.
\end{checkpoint}

\begin{lstlisting}[language=Python,caption={\texttt{check\_prime\_run.py}:
where does the $x^2 + x + 41$ prime streak break?},
label={lst:euler}]
f = lambda x: x*x + x + 41
composites = [(x, f(x)) for x in range(51) if not is_prime(f(x))]
\end{lstlisting}
The full script (in \texttt{code/check\_prime\_run.py}) defines \texttt{is\_prime}
and \texttt{least\_factor} and prints every composite. Its real output:
\begin{lstlisting}[numbers=none]
f(x) = x^2 + x + 41
x = 0..39: all prime (40 consecutive values)
x = 40: f(40) = 1681 = 41 * 41  <-- first composite
x = 41: f(41) = 1763 = 41 * 43
x = 44: f(44) = 2021 = 43 * 47
x = 49: f(49) = 2491 = 47 * 53
4 composites in x = 0..50
\end{lstlisting}

The machine confirms the algebra: $x \equiv 0$ or $-1 \pmod{41}$ forces
$41 \mid f(x)$ (that is $x = 40, 41$; the $x = 44$ and $x = 49$ failures
arrive through other prime factors). Forty examples built the conjecture;
one number killed it. That is the division of labour between examples and
proofs for the rest of this book: \textbf{examples inform, proofs establish,
counterexamples demolish.}

\section{Mathematical writing}\label{sec:writing}
\index{mathematical writing}

A proof is read by a person, and graded by one too. The rules that make it
checkable are few and mechanical:

\begin{rulebox}[Writing rules]
\index{proof!writing rules}
\begin{enumerate}
\item \textbf{Complete sentences.} A proof is English with symbols inside the
sentences, not a pile of symbols. Every line parses as a sentence.
\item \textbf{State the claim and the technique first.} ``We prove the
contrapositive.'' ``Assume, for contradiction, that\dots'' One sentence of
orientation replaces a paragraph of guessing.
\item \textbf{Name every object.} ``Let $n = 2k$ for an integer $k$.'' An
unnamed letter is an uncontrolled assumption.
\item \textbf{Displayed math for anything long.} Multi-line algebra goes
\[ \text{on its own line} \]
so the reader can verify it out of the prose flow.
\item \textbf{Ban unjustified ``clearly/obviously''.} If it is obvious, proving
it takes one line --- write the line. If it is not, you were about to lose the
reader (or the point).
\end{enumerate}
\end{rulebox}

Figure~\ref{fig:markup} applies the rules to a real flawed paragraph --- the
kind that loses marks not because the idea is wrong but because nothing in it
is licensed.

\begin{figure}[ht]\centering
% Figure: a flawed "proof" with defect marks, and what each mark means.
% Source for Figure~\ref{fig:markup}. TikZ only; regenerable by compiling main.tex.
\begin{tikzpicture}[
  bad/.style={draw,rounded corners=2pt,inner sep=7pt,font=\small,text width=52mm,align=flush left},
  ann/.style={font=\footnotesize,text=slate,text width=44mm,align=flush left},
]
\node[bad] (t) at (0,0) {%
``\textbf{Claim:} if $n^2$ is even then $n$ is even.
\textbf{Proof:} $n^2$ is even, so $n^2 = 2k$.
\defect{1}~Then $n = \sqrt{2k}$.
\defect{2}~Since $n^2$ is even, $n$ is even.
\defect{3}~So the statement holds for all $n$.
\defect{4}~Check $n = 4$: $4^2 = 16$ is even and $4$ is
even, so it works.''};
\node[ann,anchor=north west] at (3.4,0.7) {\defect{1}~No
licence for this step: $\sqrt{2k}$ need not be an integer, and the detour never
touches the parity of $n$ --- the thing the claim is about.};
\node[ann,anchor=north west] at (3.4,-0.75) {\defect{2}~Circular:
the conclusion is restated as its own justification. ``Since $n^2$ is even, $n$
is even'' is exactly the claim under proof.};
\node[ann,anchor=north west] at (3.4,-2.2) {\defect{3}~$n$ was
never fixed. The proof must open with ``let $n$ be an integer''; an unnamed
variable quantifies nothing.};
\node[ann,anchor=north west] at (3.4,-3.65) {\defect{4}~Proof by
example. One check verifies $n=4$ and says nothing about the other infinitely
many integers.};
\end{tikzpicture}

\caption{A flawed ``proof'' of the claim of Example~\ref{ex:n2even}, with four
defects marked and explained. Notice that three of the four defects are about
\emph{licence} --- steps nobody gave the writer permission to take --- and the
fourth is evidence posing as argument. The corrected proof is
Example~\ref{ex:n2even}.}
\label{fig:markup}
\end{figure}

\section{Strategy before proof}\label{sec:strategyfirst}
\index{strategy before proof}

One convention of this book deserves its own section. Every worked example here
opens with a \textbf{Strategy} paragraph --- what to try, why this technique
fits the claim's shape --- and only then the \textbf{Solution}, the cleaned
proof. The strategy is scaffolding: it is how the solution was found, and it is
the part no polished textbook ever shows you. In your own write-ups (contests,
the portfolio of Chapter~17) you submit only the proof --- but you should still
be writing the strategy somewhere, on scratch paper or in your head, because
choosing the technique \emph{is} most of the problem. Figure~\ref{fig:flowchart}
is the book's choosing procedure in one diagram.

\begin{figure}[ht]\centering
% Figure: which technique fits the goal's shape.
% Source for Figure~\ref{fig:flowchart}. TikZ only; regenerable by compiling main.tex.
\begin{tikzpicture}[
  goal/.style={draw,rounded corners=6pt,inner sep=5pt,font=\small},
  ask/.style={draw=accent,rounded corners=3pt,inner sep=3pt,font=\footnotesize,align=center},
  tech/.style={draw,rounded corners=2pt,inner sep=5pt,font=\small,text width=44mm,align=flush left},
  yes/.style={->,accent,thick},
  no/.style={->,rulegray},
  lab/.style={font=\footnotesize\sffamily,text=slate},
]
\node[goal] (start) at (0,0) {the claim you must prove};
\node[ask] (q1) at (0,-1.7) {is it $P \Rightarrow Q$?};
\node[ask] (q2) at (0,-4.3) {is it ``no object\dots''\\ or a negation?};
\node[ask] (q3) at (0,-6.9) {is it $\exists x\,P(x)$?};
\node[tech] (direct) at (5.9,-0.85) {\textbf{Direct / contrapositive.} If the
hypothesis unfolds into usable information, assume $P$ and derive $Q$; if the
\emph{conclusion} unfolds better, assume $\lnot Q$ and derive $\lnot P$.};
\node[tech] (contra) at (5.9,-3.7) {\textbf{Contradiction.} Assume the claim
fails and squeeze until it collides with a definition, an axiom, or arithmetic.};
\node[tech] (wit) at (5.9,-6.5) {\textbf{Witness.} Exhibit one $x$ and check it.
For $\exists!$, add a second act: any two witnesses must coincide.};
\node[tech] (forall) at (5.9,-9.1) {\textbf{Otherwise} ($\forall$, $\iff$, ``for
all but finitely many''): one argument covering every case, or two directions
for an iff. To \emph{refute} a $\forall$, hunt one counterexample.};
\draw[yes] (start) -- (q1);
\draw[yes] (q1.east) -- node[lab,above]{yes} (direct.west);
\draw[no] (q1) -- node[lab,right]{no} (q2);
\draw[yes] (q2.east) -- node[lab,above]{yes} (contra.west);
\draw[no] (q2) -- node[lab,right]{no} (q3);
\draw[yes] (q3.east) -- node[lab,above]{yes} (wit.west);
\draw[no] (q3.south) |- node[lab,pos=0.25,left]{no} (forall.west);
\end{tikzpicture}

\caption{Choosing a technique by the goal's shape. Notice what is \emph{not} on
the chart: ``proof by induction'' and friends arrive in Chapter~3, and every
shape still lands on a technique --- there is no ``I don't know where to start''
branch. When stuck, write the negation: the shape you see there often names the
route.}
\label{fig:flowchart}
\end{figure}

\section{Chapter summary}
\begin{table}[ht]\centering
\begin{tabularx}{\linewidth}{@{}lYYY@{}}
\toprule
\tablehead\thead{Goal shape} & \thead{Technique} & \thead{First line to write} & \thead{Killer detail}\\
\midrule
$P \Rightarrow Q$, hypothesis unfolds & direct & ``Assume $P$.'' & unfold definitions\\
$P \Rightarrow Q$, hypothesis has no handle & contrapositive & ``We prove the contrapositive: assume $\lnot Q$.'' & say the route aloud\\
``no object satisfies\dots''; $\lnot P$ & contradiction & ``Assume the claim fails.'' & find the collision\\
hypothesis splits into shapes & cases & ``Every $x$ is \dots\ or \dots; cases cover all.'' & prove exhaustiveness\\
$P \iff Q$ & two directions & ``($\Rightarrow$) \dots'' then ``($\Leftarrow$) \dots'' & both halves, each labelled\\
$\exists x\,P(x)$ & witness & ``Take $x = \dots$'' & check it works\\
$\exists! x\,P(x)$ & witness + uniqueness & ``Take $x = \dots$; now let $x_1,x_2$ both work.'' & force $x_1 = x_2$\\
refute $\forall x\,P(x)$ & counterexample & ``$x = \dots$ fails.'' & smallest candidate first\\
\bottomrule
\end{tabularx}
\caption{The technique menu on one page: the shape of the goal names the proof.
The last column is where each technique most often breaks in student work.}
\label{tab:techniques}
\end{table}

\section{Common mistakes}\label{sec:mistakes2}
\index{common mistakes}
\begin{itemize}
\item \textbf{Half a biconditional.} ``$P$ iff $Q$'' proved only in the direction
that was convenient. The other direction is a separate claim --- and usually the
one that carries the content. Label ($\Rightarrow$) and ($\Leftarrow$) so a
missing half cannot hide (Section~\ref{sec:iff}).
\item \textbf{Circular argument.} Using the goal --- or a statement that secretly
contains it --- as a step. Marked \defect{2} in
Figure~\ref{fig:markup}: ``since $n^2$ is even, $n$ is even'' asserts exactly
what was to be shown. Every line needs a licence from \emph{earlier} lines.
\item \textbf{Proof by example.} Checking $x = 0,1,2,\dots,39$ proves $x = 40$
nothing --- Figure~\ref{fig:euler}'s streak died at exactly $x = 40$. Examples
inform conjectures; they never close a universal claim.
\item \textbf{Unfounded ``clearly''.} ``Clearly $n$ is even'' is where proofs go
to hide the one step they cannot justify. Replace each ``clearly'' with its
one-line justification, or admit the step is not there.
\end{itemize}

\section{Additional exercises}
\begin{exercise}[easy]
Prove directly: for every integer $n$, if $n$ is odd then $3n + 1$ is even.
\end{exercise}
\begin{exercise}[easy]
Prove directly: the product of two odd integers is odd.
\end{exercise}
\begin{exercise}[easy]
Disprove: for every $n \ge 0$, $n^2 - n + 11$ is prime. (Find the smallest
counterexample --- it arrives sooner than Euler's did.)
\end{exercise}
\begin{exercise}[easy]
Prove $\exists\, x, y \in \Z$ with $x^2 - y^2 = 9$; then prove the solution with
$x > y > 0$ and $x + y = 9$ is unique.
\end{exercise}
\begin{exercise}[medium]
Prove by cases: for all real $x, y$, $|x + y| \le |x| + |y|$.
(Split on the signs of $x + y$, or of $x$ and $y$; AMC~12~2024~A~P6 is a
contest instance of sign-casework~\cite{aops-amc}.)
\end{exercise}
\begin{exercise}[medium]
Prove by contrapositive: for integers $a, b$, if $a + b$ is odd then exactly one
of $a, b$ is even. (Careful --- ``exactly one'' is two claims.)
\end{exercise}
\begin{exercise}[medium]
Prove there are no integers $x, y$ with $6x + 9y = 1$. Which single lens does
all the work? (AIME~2010~I~P5 applies the same one~\cite{aops-amc}.)
\end{exercise}
\begin{exercise}[medium]
Prove: for every integer $n$, $n$ is odd if and only if $n^2$ is odd.
State the technique used in each direction.
\end{exercise}
\begin{exercise}[hard]
Prove, by two different methods, that no integers $x, y$ satisfy
$x^2 - 4y = 2$: (i) modulo $4$ as in Example~\ref{ex:mod4}, and (ii) a parity
argument that never mentions residues. Then state in one sentence which method
generalises to $x^2 - 9y = 3$ and why.
\end{exercise}
\begin{exercise}[hard]
Extend \texttt{code/check\_prime\_run.py} to test other prime-producing
polynomials: $g(x) = x^2 + x + 17$ on $x = 0..40$ and $h(x) = 2x^2 + 29$ on
$x = 0..60$. Report each first composite and explain --- by factoring like
$f(40) = 41^2$ --- why it was forced.
\end{exercise}
\begin{exercise}[hard]
Euclid's argument. Assume there are only finitely many primes
$p_1, p_2, \dots, p_k$, and let $N = p_1 p_2 \cdots p_k + 1$. Prove that $N$
has a prime factor not among $p_1,\dots,p_k$, conclude infinitely many primes
exist, and explain why ``$N$ is prime'' is \emph{not} the right claim to make.
\end{exercise}
\begin{exercise}[easy]
For each claim, name the proof technique the shape suggests (no proof needed):
(i) ``if $3n+2$ is odd then $n$ is odd''; (ii) ``no integer is both even and
odd''; (iii) ``$n$ is even iff $n{+}2$ is even''; (iv) ``for every nonzero
digit $d$, $d^2$'s last digit is in $\{1,4,5,6,9\}$''.
\end{exercise}
\begin{exercise}[medium]
Prove by cases on $n \bmod 3$: for every integer $n$, $n^3 - n$ is divisible
by $3$. Write the case-exhaustiveness line explicitly.
\end{exercise}
\begin{exercise}[medium]
Disprove or prove, choosing the right tool: ``for every prime $p$, $2^p - 1$
is prime.'' If false, give the smallest counterexample (check
$p = 2,3,5,7,11$).
\end{exercise}
\begin{exercise}[hard]
Prove: for every integer $n$, $n$ is odd iff $n^2 \equiv 1 \pmod{8}$.
Label both directions and say which one needs cases on $n = 4k{+}1$ vs
$4k{+}3$.
\end{exercise}

\section*{Further reading}
\begin{itemize}
\item R.~Hammack, \emph{Book of Proof}, 3rd ed., Ch.~4--10 ---
\url{https://richardhammack.github.io/BookOfProof/} --- free PDF; the
technique-by-technique spine this chapter mirrors, with many more worked
proofs.
\item E.~Lehman, F.~T.~Leighton, A.~R.~Meyer, \emph{Mathematics for Computer
Science}, Ch.~1--2 ---
\url{https://courses.csail.mit.edu/6.042/spring18/mcs.pdf} --- free PDF;
``patterns of proof'' done CS-style, fast and rigorous.
\item E.~Chen, \emph{Remarks on English} ---
\url{https://web.evanchen.cc/handouts/english/english.pdf} --- the
write-up standards handout; read it once per section of this chapter.
\item Stanford CS103, \emph{Mathematical Foundations of Computing} ---
\url{https://web.stanford.edu/class/cs103/} --- free reader; its
proof-writing guides cover exactly this chapter's ground.
\item MIT OpenCourseWare, \emph{6.042J Mathematics for Computer Science}
(Fall 2010) --- \url{https://ocw.mit.edu/6-042JF10} --- lecture notes and
problem sets for the proofs-first half of the course.
\item AoPS, \emph{How to Write a Solution} --- contest write-up advice;
pairs with \S\ref{sec:writing}.
\url{https://artofproblemsolving.com/articles/how-to-write-solution}
\end{itemize}

\begin{rulebox}[Watch instead]
\begin{itemize}
\item \textbf{TrevTutor --- ``[Discrete Mathematics] Direct Proofs
Examples''}, 6:59, \url{https://www.youtube.com/watch?v=uDJfx4bK3Jc} ---
covers \S\ref{sec:direct}: two worked direct proofs including divisibility.
(Verified 2026-10-11.)
\item \textbf{TrevTutor --- ``Can You Prove? Proofs in Discrete Math and
Real Analysis''} (playlist),
\url{https://www.youtube.com/playlist?list=PLDDGPdw7e6AjLVCdoobmqvyZZpthHGBJ6}
--- covers \S\ref{sec:cases}--\ref{sec:existence}: contraposition,
contradiction, cases and existence proofs in short segments.
(Verified 2026-10-11.)
\end{itemize}
\end{rulebox}

\chapternotes
The technique-per-chapter architecture (direct, contrapositive, contradiction,
then the non-conditional shapes) follows Hammack's proof chapters, whose
``Mathematical Writing'' section inspired Section~\ref{sec:writing} and whose
disproof chapter put counterexamples on equal footing with
proofs~\cite{hammack}. Velleman's third chapter organises the same material by
the logical form of the goal --- our Figure~\ref{fig:flowchart} is a compressed
version of his summary tables~\cite{velleman}. MIT 6.042J's opening chapters
supplied the ``axioms vs.\ proven tools'' framing and the
squares-mod-$4$ lens~\cite{mcs}. The writing rules of
Section~\ref{sec:writing} and the strategy/solution split of
Section~\ref{sec:strategyfirst} follow Chen's \emph{Remarks on English}, whose
side-by-side marked-up solutions modelled Figure~\ref{fig:markup}~\cite{chen-english},
and the mistakes list in OTIS Excerpts ch.~1~\cite{otis}; Cut-the-Knot's
fallacy pages supply the gallery this section samples~\cite{cuttheknot}.
Euler's polynomial example is standard lore (Euler, 1772); problem citations
(AMC~12~2024~B~P14, AMC~12~2023~A~P3, AIME~2010~I~P5, AMC~12~2024~A~P6) point to
the AoPS archive~\cite{aops-amc}.
